\section{The two-puncture observable and its finite spin formula}\label{c7:finite}

We construct the finite cylinder observable that determines the nesting
partition. Orienting the loops expresses it as a pairing of two vacuum
vectors. We then identify the scalar factor that clears the denominators
of these particular vectors and use their polynomial degrees and fusion
values to determine the spin formula. The same vector statement will
also supply the regularity needed for the two marked bonds.

Keep the honeycomb cylinder of circumference $M$, first with finite height on either side of a cut along a horizontal tessellation line. Mark two face centers (tessellation vertices) $O,O'$ on the cut, the first at the seam, dividing the sites into consecutive groups $A,B$ in eastward order starting at $O$. Close both rims. Use any collection of disjoint simple cycles, \emph{each separating} $O,O'$, with weight $x^{|\gamma|}g$ per cycle, $x$ the critical weight above. Write $G_{A,B}(g)$ for this partition including the empty term, with the height sent to infinity. We use $g=h+h^{-1}$ (in this part $h$ denotes a twist parameter, not the sine constant in the weights). The degree in $g$ is bounded by $M$ since the cycles must intersect the arc of the cut from $O$ to $O'$.

Here and below matrices of oriented weights and their spectral parameter $\phi$ refer to $R(\phi)$ above, $\lambda=e^{-2i\alpha}$ and $\alpha=\pi/8$. The rhombic Yang--Baxter setting and periodic dilute-loop spectra have the precedents \cite{GlazmanManolescu2017,ZhouBatchelor1997}. The current calculation of \cite[Section~6]{FeherNienhuis2018} is a related recursion under a rapidity-symmetry hypothesis; our finite observable and its symmetry are derived below.

\subsection{The oriented observable and its vacuum limits}

Put $q=e^{2i\alpha}=\lambda^{-1}$, introduce site variables $x_j=e^{2i\phi_j}$ (not vertex fugacities), and, with $w=e^{2i\xi}$, set
\[
T_p(w)=\mathrm{Tr}_a\, p^{s_a} R_{a1}(\xi-\phi_1)R_{a2}(\xi-\phi_2)\cdots R_{aM}(\xi-\phi_M),
 \qquad p\in\{i,-i\}.
\]
Here $s_j$ is the diagonal spin, index $a$ the auxiliary space. We work
on the total site charge zero subspace
\[
\mathcal H_M=\operatorname{Span}\{|\sigma\rangle:
\sigma\in\{-1,0,1\}^M,\ \sum_j\sigma_j=0\}
\]
and denote its empty vector by $|0\rangle$. Choices of square roots of
the variables are understood. For each fixed $p$, the rows commute by
Lemma~\ref{j:local-exchange-fusion}. The vacuum matrix element satisfies
\[
\langle0|T_{p_1}(w_1)\cdots T_{p_L}(w_L)|0\rangle=1
\]
for arbitrary words in the two twists, at regular parameters. Indeed the
oriented phases in the blocks can still be drawn on the reference tiling
with fixed angle $\pi/3$, keeping the angle-independent phases after
removing the telescoping factor and varying only the scalar weights of
the tile pictures. Each contractible loop has phase $\pm i$ before seam
signs, and each noncontractible one has turn zero and winding $\pm1$.
Seam factors for either orientation are those for uniform twist $i$,
times the same sign from the layers of twist $-i$. All loops therefore
cancel.

We will use this word identity in two ways. First, every linear
combination $A$ of operators $T_p(w)-I$, even with the two twists mixed,
has zero as an eigenvalue: $\langle0|A^k|0\rangle=0$ for $k\ge1$, and
the characteristic polynomial would contradict this if its constant
term were nonzero. Second, fix one twist $p$. Choose finitely many row
deviations spanning the linear span of all $T_p(w)-I$ in
$\operatorname{End}\mathcal H_M$ and simultaneously triangularize them. For every
choice of coefficients, one diagonal character of the corresponding
linear combination vanishes. A finite union of proper hyperplanes
cannot cover the coefficient space, so one character vanishes
identically. Thus this fixed-twist family has the simultaneous
eigencharacter $T_p(w)=1$.

For physical site variables all one, use twist $p=-i$ above the cut and $p=i$ below, with the insertion
\[
D_A=(ih)^{\sum_{j\in A}s_j}.
\]
To check the weight, orient a contractible simple cycle counterclockwise or a noncontractible one eastward, and let $I_1,I_2$ indicate whether $O,O'$ belong to its left component (interior or upper, respectively). The signed cut crossings in $A$ (downward positive) give $I_2-I_1$. Signed seam crossings above and below $O$ (eastward positive) give $-I_1,I_1$ in the contractible case and $1-I_1,I_1$ otherwise. These numbers follow by moving along the corresponding transverse arcs from component to component. Including turn, seam and insertion, the phase is $(-i)(-1)^{I_1}(ih)^{I_2-I_1}$, with its inverse for reversed orientation. Thus it cancels unless the cycle separates the marks, and then sums to $g$.

The top and bottom vacuum vectors tend to finite limits
$L_-^{\rm phys},R_+^{\rm phys}$; the subscripts $+,-$ on vacuum vectors
stand for twists $i,-i$. Untouched closed loops separately cancel,
leaving disjoint arcs connecting ports on the cut in each half-cylinder.
Any lift of such an arc has port separation less than $M$: otherwise its
endpoints interlace those of a period translate, contradicting
noncrossing. Products of finitely many half-plane arch sums $H(d)$ from
Section~\ref{j:local} therefore dominate the absolute weights.

These physical vectors have the simultaneous eigencharacter $T=1$ for
their respective twists. For example, insert a regular row $T_i(w)$
just below the cut and commute it to the far closed rim. With the rim
ports empty that row uses no double turn, so its absolute weights are
bounded by the physical weights times a constant depending on $M,w$.
An arc affected by the row must reach the receding rim. Its total weight
tends to zero by the same arch summability. The inserted row thus leaves
the limit unchanged, and the top vector is treated in the same way.
Both vacuum components are 1. This proves the height limit and finite
coefficient sums (for example by taking $g>0$), with
\[
G_{A,B}=L_-^{\rm phys}D_A R_+^{\rm phys}.
\]

\subsection{Fusion and the algebraic vacuum projector}

Write $H_p=T_p-I$. We first record the behavior of a local matrix when
$z=w/x_j$ tends to zero or infinity. Directly from the sine weights,
\[
\begin{array}{c|c|c}
 & (a,b,v)&(u,d)\\ \hline
z\to0&(q^{-3},q^3,1)&O(z^{1/2})\\
z\to\infty&(q^3,q^{-3},1)&O(z^{-1/2}).
\end{array}
\]
The entries are holomorphic in $z^{1/2}$ at zero and in $z^{-1/2}$ at
infinity. On a state with auxiliary spin $r$ and site spin $s$, the
diagonal limit is $q^{3rs}$ at infinity and $q^{-3rs}$ at zero. The
only surviving off-diagonal entry at infinity changes the auxiliary
spin from $-1$ to $1$ and the site spin from $1$ to $-1$ (row minus
column); at zero this move is reversed. In particular every surviving
move at an inactive site preserves its vacancy parity.

Square-root factors occur only in the single turns, which change
auxiliary vacancy parity. Their number in an auxiliary trace is even,
so $H_p(w)$ is rational in $w$, with possible simple poles
$q^5x_j,q^6x_j$ at generic sites. The same observation for a site says
that changing a square-root sign conjugates $T_p$ by
$P_j=(-1)^{s_j^2}$. At either endpoint of the row parameter, a closed
auxiliary trace cannot use a strictly one-directional off-diagonal
move. On $\mathcal H_M$ the remaining diagonal trace is
$\sum_{r=-1}^1p^r q^{\pm3r\sum_j s_j}=1$. Thus
$T_p(0)=T_p(\infty)=I$. With $C_\ell(z,b)=z/(z-q^\ell b)$ we obtain
\[
H_p(w)=\sum_{j=1}^M\sum_{t=0}^1 C_{5+t}(w,x_j) U_{jt},
\]
where the coefficient operators commute. The basis functions vanish at
zero and have distinct poles; the condition at infinity additionally
gives $\sum_{j,t}U_{jt}=0$.

The two embeddings and row fusion above give at any site $x_j$, generically,
\[
T_p(x_j)T_p(q^2x_j)=T_p(qx_j),\qquad T_p(x_j)T_p(q^3x_j)=I.
\]
Indeed the combined auxiliary monodromy has local factors
$R_{2k}(\phi)R_{1k}(\phi+\delta\alpha)$, with the lower parameter on
the left. Lemma~\ref{j:local-restriction} identifies their restriction
to $\operatorname{im}J_\delta$ and shows that the quotient factor at
site $j$ is zero: its full image is already in
$\operatorname{im}J_\delta\otimes V_j$. The quotient monodromy therefore
has zero trace. On the restriction the twist is the fused twist, or the
identity for $\delta=3$, giving the two formulas. This is a trace of
the restriction transported through $J_\delta$, so no isometry factor
enters. Square roots can be chosen consistently; shifting a whole row
parameter by $\pi$ leaves $T_p$ unchanged.

Set $T_j=T_p(x_j)$. Expanding the fusion formulas in the commuting
coefficients $U_{kt}$ gives a homogeneous system $\mathcal L(T)U=0$.
Its two rows for site $j$, with column site $k$, type $t$, and
$C_\ell=C_\ell(x_j,x_k)$, are
\[
\begin{split}
 C_{5+t}-C_{4+t}+T_j C_{3+t},\\
 C_{4+t}+T_j(C_{2+t}-C_{3+t}).
 \end{split}
\]
The second is the rank-one fusion equation minus the first. At $T=1$
these rows are $C_5+C_3-C_4$ on the same row and column list
$Y=(x_1,q x_1,\ldots,x_M,q x_M)$. Subtracting the constant $1/2$ gives
the skew kernel
\[
S(u,v)=\frac{(u-v)(2(u+v)^2-d(u,v))}{2(u+v)d(u,v)},\qquad d(u,v)=u^2+\sqrt2uv+v^2.
\]
The adjugate of an even skew matrix is skew. Adding a matrix with
constant entries therefore does not change its determinant, and
$\det\mathcal L(1)=(\operatorname{Pf}S(Y))^2$. Whenever this scalar
denominator is nonzero,
\[
\Pi_p=\frac{\det\mathcal L(T)}{\det\mathcal L(1)}
\]
is a projector onto the common fixed space. Indeed the determinant is
taken in the commutative algebra generated by the rows. The adjugate
identity for $\mathcal L(T)U=0$ shows that $\Pi_p$ annihilates every
$U_{jt}$, while on their joint kernel $T_j=1$ and $\Pi_p=1$. Hence its
image is exactly that kernel and it is idempotent. Evaluating its
polynomial numerator by the arbitrary-word vacuum identity also gives
\[
\langle0|\Pi_p|0\rangle
 =\langle0|\Pi_p\Pi_{p'}|0\rangle=1
\qquad(p,p'\in\{i,-i\}).
\]

\subsection{The scalar normalizers}

The denominator of the projector is governed by three scalar functions.
The factorization below separates the factor $N$, which will clear the
two selected vacuum vectors, from the additional factor $N_2^*$ present
for the full projector.

We use the Pfaffian of an even skew matrix: it is the signed matching
sum, equivalently the coefficient of the oriented volume form in the
normalized top power of its two-form. It transforms under simultaneous
row and column change by the determinant, and its square is the
determinant, as follows by reduction to two-dimensional skew blocks.
For the site list $X$, put
$\mathcal P(X)=\prod_{j<k}d(x_j,x_k)/(x_j-x_k)$ and define
\[
\begin{split}
 N(X)&=\mathcal P(X)\operatorname{Pf}_{j,k}\frac{x_j^2-x_k^2}{d(x_j,x_k)},\\
 N_2^*(X)&=\mathcal P(X)^2\det_{j,k}\frac{(2+\sqrt2)x_jx_k}{d(x_j,x_k)},\\
 \mathcal J(X)&=S(1,q)^{-M}\operatorname{Pf} S(Y)\prod_{j<k}P_*(x_j,x_k),\\
 P_*(u,v)&=\frac{(u+v)(u^2+v^2)d(u,v)^2}{(u-v)^2(u^2-\sqrt2uv+v^2)} .
 \end{split}
\]
When $M$ is odd, pad the Pfaffian defining $N$ by a last index whose
column entries above the diagonal are 1. All empty-list normalizations
are 1. Write $N_1=N$, $N_2=N_2^*$, $N_3=\mathcal J$, and set
\[
\ell_2(u,z)=z-q^5u,\qquad
\ell_3(u,z)=(z-q^5u)(z-q^6u).
\]
\begin{lemma}[Scalar normalizers]\label{j:scalar-normalizers}
For every list size $M$, the functions $N_l$ are symmetric homogeneous
polynomials of individual degree at most $l(M-1)$ for $M\ge1$, and
\[
\mathcal J=N N_2^*.
\]
For $l=1,2$ and $X=(u,Z)$ of size $M\ge1$, their deletion values are
\begin{equation}\label{j:normalizer-deletion}
N_l(0,Z)=\left(\prod_{z\in Z}z^l\right)N_l(Z),
\qquad [u^{l(M-1)}]N_l(u,Z)=N_l(Z).
\end{equation}
For $\delta=2,3$ all three satisfy
\begin{equation}\label{j:normalizer-fusion}
N_l(u,q^\delta u,Z)=N_l(u,q^\delta u)\ \prod_{z\in Z}\ell_\delta(u,z)^l\ N_l(X_{\rm red}),\qquad
 X_{\rm red}=\begin{cases}(qu,Z)&\delta=2,\\Z&\delta=3.\end{cases}
\end{equation}
At two sites,
\[
N(u,v)=u+v,\qquad
N_2^*(u,v)=u^2+2(1+\sqrt2)uv+v^2.
\]
Finally, $N$ and $N_2^*$ are nonzero at every homogeneous list
$(b,\ldots,b)$ with $b\ne0$.
\end{lemma}
\begin{proof}
We first remove the apparent poles. For $N$ and $N_2^*$ all $d$
denominators are cleared by $\mathcal P$ or $\mathcal P^2$. The
remaining numerator vanishes on equal sites by alternating rows, or by
alternating rows and columns, so it is divisible by the corresponding
Vandermonde factors.
The sign from a site interchange in the Pfaffian is canceled by that
of $\mathcal P$; the determinant has the corresponding row and column
signs. For $\mathcal J$, interchanging two paired blocks is an even
permutation. Thus all three functions are symmetric.

For $\mathcal J$, fix a pair of sites $u,v$ with $v=q^\delta u$ and
keep all other sites generic. In the skew kernel on the paired list
$Y$, repeated rows and columns give a zero of order at least two at
$\delta=0$ and at least one at $\delta=\pm1$. At $\delta=2$ the only singular
cross coupling is between $u,qv$, so there is at most a simple
Pfaffian pole. At $\delta=3$ there is at most a double pole, and the
reflected cases have the same bounds. At $\delta=4$ the $2\times2$
cross residue, up to a common scalar, has diagonal entries 1 and
off-diagonal entries $-1$. Its rank is one, so the potential double
pole cancels. The factors in
\[
P_*(u,v)=
\frac{(v-q^4u)(v-q^2u)(v-q^6u)
      (v-q^3u)^2(v-q^5u)^2}
     {(v-u)^2(v-qu)(v-q^7u)}
\]
clear these possible orders. These calculations are at generic points
of each possible divisor. Unique factorization then removes every
irreducible denominator; simultaneous intersections are included by
this polynomial conclusion.

Homogeneity follows from the displayed formulas, and the individual
degree bounds follow by taking a variable to infinity. For the paired rows of $\operatorname{Pf}S(Y)$
the cross rows have identical limits, so they decouple from the rest
of the Pfaffian. The deletion rules for $N,N_2^*$ follow in the same
way at zero and infinity. The extra determinant index decouples.
For the Pfaffian, place the varied variable first. Its row on the other
sites has constant limit $-1$ at zero and $1$ at infinity. Without
padding it becomes the padding row after reordering; with padding the
two constant rows pair off since their mutual coupling is 1. The signs
cancel those from $\mathcal P$, giving \eqref{j:normalizer-deletion}.

We next prove the reductions. For $\mathcal J$ the leading pole at
$\delta=3$ removes both pairs of site indices, while at $\delta=2$ the
pole between $u,q^3u$ leaves $qu,q^2u$, the pair for the reduced site.
The internal scalar is the two-site value. The remaining factors are
\[
P_*(u,z)P_*(q^3u,z)=\ell_3(u,z)^3,\qquad
\frac{P_*(u,z)P_*(q^2u,z)}{P_*(qu,z)}=\ell_2(u,z)^3.
\]
For $\delta=3$, the $d$ pole and $\mathcal P$ give the same reduction
for the first two polynomials.

For $N$ at $\delta=2$, first divide out
$\prod_{z\in Z}(z-q^5u)$. At each $z=q^5u$ the prefactor has two
zeros, whereas the Pfaffian has at most one pole because the two
singular entries share the index $z$. After division, compare the two
sides as polynomials of degree at most $M-2$ in one spectator $z$,
inductively in $M$. The deletion rules give their values at zero and
their top coefficients. At $z=q^{\pm3}y$ for every other spectator
$y$, the already proved $\delta=3$ reduction and induction give
equality, using
\[
\frac{\ell_3(y,u)\ell_3(y,q^2u)}
     {(y-q^5u)(q^3y-q^5u)}=\ell_3(y,qu).
\]
These evaluations suffice for $M\ge3$; for $M=2$ the formula is
tautological.

For $N_2^*$ at $\delta=2$, write its determinant kernel as
\[
K(x,y)=(2+\sqrt2)\frac{\langle f(x), f(y)\rangle}{x^4+y^4},\qquad
 f(x)=(x,x^2,x^3),\quad \langle a,b\rangle=a_1b_3+a_3b_1-\sqrt2a_2b_2.
\]
The two vectors at $u,q^2u$ span the orthogonal complement of the
nonisotropic vector at $qu$. Here $u^4=(q^2u)^4=U$, while
$(qu)^4=-U$. Eliminate the first two rows and columns by a Schur
complement. On their span, the subtracted scalar factor is
$2U/[(x^4+U)(y^4+U)]$, and
\[
\frac1{x^4+y^4}-\frac{2U}{(x^4+U)(y^4+U)}
 =\frac{r(x)r(y)}{x^4+y^4},\qquad
r(x)=\frac{x^4-U}{x^4+U}.
\]
The orthogonal line is unchanged. Eliminating instead the single
$qu$ in the reduced determinant leaves the complement unchanged and
multiplies the line pairing by $1/(r(x)r(y))$. Thus multiplying each
spectator row and column of that Schur matrix by $r$ gives the first
Schur matrix. Their determinants differ by $\prod r^2$, and inserting
$\mathcal P^2$ gives \eqref{j:normalizer-fusion}.

The stated two-site formulas follow directly from the Pfaffian and
determinant. Their product agrees with $\mathcal J$: by symmetry and
homogeneity it suffices to expand at $(1,v)$ through first order.
Here
$P_*(1,v)=1+(3+3\sqrt2)v+\cdots$,
$S(1,v)=1/2+(1-\sqrt2)v+\cdots$, and
$S(1,q)=-i(2-\sqrt2)/2$. The cross contractions
$-S(1,v)S(q,qv)+S(1,qv)S(q,v)$ shift the linear coefficient by
$-\sqrt2$, giving the product. The one-site values are 1.
Induction using the common reductions now proves
$\mathcal J=NN_2^*$: in one variable there are $4(M-1)$ evaluations
at $q^{\pm2},q^{\pm3}$ times the other variables, more than its degree
$3(M-1)$.

Finally consider confluence at a nonzero homogeneous list. In log
variables the Pfaffian kernel is
$B(t)=\sinh t/(\cosh t+\cos\beta)$, with $\beta=\pi/4$, and the
determinant kernel is a positive multiple of
$(\cosh t+\cos\beta)^{-1}$. Dividing the Vandermondes takes confluent
row and column derivatives, up to nonzero constants.
Lemma~\ref{j:fourier-kernel} gives the latter kernel a strictly
positive Fourier density with all moments finite, so its derivative
matrix is a nondegenerate Gram matrix. Applying
$\cos\beta+\sin\beta\,\partial_\beta$ to that transform gives
\[
\widehat{B'}(k)=\frac{2\pi k\cosh(\beta k)}{\sinh(\pi k)}.
\]
This is also a strictly positive even density. For even $M$, the
cross matrix of even and odd derivative orders of $B$ is, up to
row and column signs, its moment Gram matrix on even powers; the
other entries vanish. For odd $M$ the padding pairs with derivative
order zero and the remaining cross matrix has the same form with the
positive weight multiplied by $k^2$. Both are nondegenerate.
\end{proof}

\subsection{The fixed vacuum lines}

We use widely separated site magnitudes to identify the fixed character
and to locate the possible poles of its projection. The following
description also records the common similarity needed when the sites of
a scale are interspersed in the cyclic order.

\begin{lemma}[Separated scales]\label{j:separated-scales}
Let positive scales $\rho_1,\ldots,\rho_r$ satisfy
$\rho_{k+1}/\rho_k\to\infty$, and suppose
$x_i/\rho_{k_i}$ tends to a finite nonzero limit. Choose consistent
local square roots. Probe with $w/\rho_j$ tending to a finite nonzero
limit for which the active local matrices are regular. The limiting
row matrices are triangular in the grading $\sum_i k_i s_i$, with row
grade at least column grade. Their diagonal grade blocks preserve every
scale total $S_k=\sum_{i:k_i=k}s_i$ and act on the active scale $j$
with effective twist
\[
p_j=p q^{3\sum_{k\ne j}\operatorname{sgn}(j-k) S_k} .
\]
Each block entry is the corresponding active-transfer entry multiplied
by $F(s^{\rm row})/F(s^{\rm col})$, using the same diagonal matrix
for every probe and twist, where
\[
F(s)=q^{(3/2)\sum_{l<i}\operatorname{sgn}(k_l-k_i)s_l s_i} ,
\]
restricted to the block in question.
\end{lemma}
\begin{proof}
Let $\Delta s_i=s_i^{\rm row}-s_i^{\rm col}$. At an inactive site on a
smaller scale than $j$, the local limit gives $\Delta s_i\le0$; on a
larger scale it gives $\Delta s_i\ge0$. Total charge conservation gives
\[
\Delta\!\left(\sum_i k_i s_i\right)
 =\sum_i(k_i-j)\Delta s_i\ge0.
\]
The inequality is strict if any inactive charge changes. Thus a
diagonal grade block fixes the inactive spins and conserves the total
on the active scale.

At an inactive site the diagonal factor is
$q^{3\operatorname{sgn}(j-k_i)h_i s_i}$, where the incoming horizontal
spin is
$h_i=h_0+\sum_{l<i}(s_l^{\rm row}-s_l^{\rm col})$.
The $h_0$ contribution changes the seam twist to $p_j$. For the
remaining contribution all changes lie on the active scale and sum
to zero. In the exponent of $F(s^{\rm row})/F(s^{\rm col})$, an
inactive site receives equal contributions from the active sites
before and after it, because those two charge changes have opposite
sums. Their combined exponent is
$3\operatorname{sgn}(j-k_i)s_i\sum_{l<i}\Delta s_l$, exactly the
remaining local factor. This proves the stated common similarity.
\end{proof}

Take all scales to be singletons and probe at each singleton's swap
point. On a diagonal block its eigenvalue is $p_j^{S_j}$. If
$S_j=\pm1$, total charge zero implies
$\sum_{k\ne j}\operatorname{sgn}(j-k)S_k$ is odd. Since
$p=q^{2}$ or $q^{6}$, the exponent of $q$ in $p_j$ is odd, and
$p_j^{S_j}\ne1$. Only the all-empty block can have the simultaneous
character, with multiplicity one. A generic linear combination of
these finitely many probe deviations therefore has a simple zero
eigenvalue in the limit and for sufficiently separated actual sites;
zero is always an eigenvalue by the vacuum argument. Rationality in
square-root coordinates proves generic simplicity, with no Jordan
extension at the fixed character.

Consequently the generic projector has rank one. Set
\[
L_p=\langle0|\Pi_p,\qquad R_p=\Pi_p|0\rangle.
\]
Vacuum normalization gives
$\Pi_p=R_pL_p$ and $L_pR_{p'}=1$ for both twists.
We use $R_+=R_i$, $L_-=L_{-i}$ and the corresponding shorthand for
the other two vectors.

\paragraph{Homogeneous sites.}
At a list $(b,\ldots,b)$ with $b\ne0$, the determinant projector has
a confluent form. The pole basis for $H_p$ consists of
$C_5(w,b),C_6(w,b)$ and their $b$-derivatives through order $M-1$;
these are independent partial fractions with value zero at $w=0$.
Use identical rapidity branches near the common site. By
Lemma~\ref{j:local-restriction}, each quotient factor in the fusion
monodromy is analytic and divisible by $w-b$. The product of the
$M$ factors is therefore $O((w-b)^M)$. Both row fusion identities,
viewed as functions of the probe $w$, hold through derivatives of
order $M-1$ at $b$.

Apply these derivatives to the rows of $\mathcal L$, differentiating
$T_p(w)$ as well, and apply the corresponding derivatives to its pole
columns. The resulting system annihilates the confluent pole
coefficients. Its scalar determinant is the limit of
$\det\mathcal L(1)$ after division by two row and two column
Vandermondes. To see its value, put
$\Delta(X)=\prod_{j<k}(x_j-x_k)$. The product
$P_*(u,v)(u-v)^2$ is analytic and nonzero near $u=v=b\ne0$.
Thus $\det\mathcal L(1)/\Delta(X)^4$ is a nonvanishing analytic
factor times $(NN_2^*)^2$, and is nonzero at the homogeneous list
by Lemma~\ref{j:scalar-normalizers}.
The same row and column divided differences in the numerator show
that the determinant quotient is the ordinary limit of the generic
$\Pi_p$: all functions involved are jointly analytic near the
differentiation points. The limit is a rank-one projector with vacuum
normalization. At $X=(1,\ldots,1)$ its fixed vectors are therefore
the physical limits $L_-^{\rm phys},R_+^{\rm phys}$ constructed above.

\subsection{Selected vacuum vectors and polynomial degrees}

The projector denominator contains both $N$ and $N_2^*$. We now show
that $N$ alone removes the poles of the particular column $R_i$ and
row $L_{-i}$, after accounting for their square-root parity. This
component statement will give the degree bound for the scalar pairing.

\begin{lemma}[Selected vacuum vectors]\label{j:selected-vacuum-vectors}
Let $M\ge1$, let $X=(x_1,\ldots,x_M)$ be generic and nonzero, and choose
$t_j$ with $t_j^2=x_j$. For
$\sigma\in\{-1,0,1\}^M$ with $\sum_j\sigma_j=0$, define from the
vacuum-normalized vectors $R_i=\Pi_i|0\rangle$ and
$L_{-i}=\langle0|\Pi_{-i}$ the functions
\[
V_\sigma^R(X)=
\frac{N(X)(R_i)_\sigma}{\prod_jt_j^{\sigma_j^2}},
\qquad
V_\sigma^L(X)=
\frac{N(X)(L_{-i})_\sigma}{\prod_jt_j^{\sigma_j^2}}.
\]
They are independent of the choices of square-root signs and extend
to polynomials in $\mathbb C[x_1,\ldots,x_M]$ satisfying
\[
\deg_{x_j}V_\sigma^\epsilon\le M-1-\sigma_j^2
\qquad(\epsilon=R,L).
\]
Write $X_{\widehat j}$ and $\widehat\sigma$ for deletion of coordinate
$j$, and set $V_\varnothing^R=V_\varnothing^L=1$. If $\sigma_j=0$,
then for either superscript
\[
\left.V_\sigma(X)\right|_{x_j=0}
 =\left(\prod_{k\ne j}x_k\right)V_{\widehat\sigma}(X_{\widehat j}),
\qquad
[x_j^{M-1}]V_\sigma(X)=V_{\widehat\sigma}(X_{\widehat j}).
\]
\end{lemma}
\begin{proof}
Site sign covariance gives
$\Pi_p(-t_j)=P_j\Pi_p(t_j)P_j$. Since the empty vector has even
vacancy parity,
\[
(R_p)_\sigma(-t_j)=(-1)^{\sigma_j^2}(R_p)_\sigma(t_j),\qquad
(L_p)_\sigma(-t_j)=(-1)^{\sigma_j^2}(L_p)_\sigma(t_j).
\]
The determinant formula makes the functions in the statement rational
in the $t_j$, and they are invariant under every sign change. The
invariant rational field is
$\mathbb C(t_1^2,\ldots,t_M^2)=\mathbb C(X)$, so they are rational in
the ordinary variables.

We locate their possible nonzero finite poles one variable at a time,
with generic spectators. Fix all sites except $\zeta$ on widely
separated singleton scales. At a putative pole, suppose a regular
linear combination of commuting probe deviations has a simple zero
eigenvalue. Its eigenprojection agrees nearby with $\Pi_p$, since it
projects onto the same fixed line, and is regular there. Such a probe
rules out a pole of every normalized row and column.

If the varied site shares a scale with one spectator, the singleton
tests in Lemma~\ref{j:separated-scales} require every other site to
be empty in a diagonal block carrying the fixed character. The active
pair then has total zero and original twist $p$. Write its coordinates
in left-to-right order as $u,v$ and put $z=u/v$. The probe $w=u$ is
the twisted swap reduction
$R(\phi_u-\phi_v)\operatorname{diag}(p^{s_{\rm first}})$. Its neutral
block has eigenvalues $1,a,-a$, where
\[
a=q^3\frac{(z-q^2)(z-q^3)}{(z-q^6)(z-q^5)}.
\]
Indeed it has eigenvalue 1 by the vacuum character, trace 1, and
trace of its square $1+2a^2$. In this pair block, normalize the
right and left fixed vectors $r,l$ at the state $(0,0)$, and let their
subscripts $+,-$ denote states $(1,-1),(-1,1)$. Direct multiplication
gives
\[
r_-=ir_+,\qquad l_-=il_+,\qquad
r_+=\frac{d_{\rm tile}\lambda}{1-ia/p},\qquad
l_+=\frac{d_{\rm tile}p/\lambda}{1-iap}.
\]
Here $d_{\rm tile}$ is the local single-turn weight, a nonzero constant
times $\sqrt z(z-1)/[(z-q^6)(z-q^5)]$, and is distinct from the
quadratic $d(u,v)$. These formulas also give $lr=1$.

At a pole of $a$, the reverse probe $w=v$ is regular and has
$a(1/z)=0$; moving the seam cyclically only conjugates the pair.
Thus the fixed eigenvalue is simple there. If
$D(z)=(z-q^5)(z-q^6)$, the remaining crossings are explicit:
\[
a-1=\frac{(q^3-1)(z^2-1)}{D(z)},\qquad
a+1=\frac{(q^3+1)(z^2+2(1+\sqrt2)z+1)}{D(z)}.
\]
At $z=1$ use identical rapidity branches and also probe at $w=qu$.
On the two-dimensional 1-space of the twisted swap, this probe has
trace $1+v_{\rm tile}^2+a_{\rm tile}^2+b_{\rm tile}^2$, with weights
at $\alpha$. Its polar diagonal entries are $v_{\rm tile}^2$ and
its polar cross entries are $(b_{\rm tile}^2+a_{\rm tile}^2)p^{\pm1}$.
One eigenvalue is 1 and the other is
$v_{\rm tile}^2+a_{\rm tile}^2+b_{\rm tile}^2\ne1$. Joint simplicity
therefore holds. All these pair calculations concern the diagonal
block with empty spectators.

To make the localization uniform, consider any sequence of generic
spectator lists whose scale gaps tend to infinity and any associated
putative poles $\zeta\ne0,\infty$. After a subsequence, $\zeta$ is
either on a separated singleton scale or in a pair whose ratio tends
to a finite nonzero limit. The corresponding scale probes converge
locally analytically in a square-root chart. In the singleton case
and at every jointly simple pair ratio, a probe combination has a
simple zero in the limit and hence for sufficiently large gaps.
The only ratios left are $-1$ and the two negative roots of
$z^2+2(1+\sqrt2)z+1$.

Around one of these three ratios take a small complex $z$-disk and
use the pair probe $w=u$ together with the spectator singleton probes.
For generic constants set
\[
A(z)=\sum_\nu c_\nu(T_p(w_\nu)-I),\qquad
d_A(z)=\left.\partial_\rho\det(\rho I-A(z))\right|_{\rho=0}.
\]
Here $d_A$ is a scalar derivative, unrelated to the insertion $D_A$.
The singleton tests keep every block with a charged spectator
invertible for $A$. In the remaining block the singleton probes are
identity and the pair eigenvalues cross transversely by the displayed
factorizations. Thus the limiting $d_A$ has a simple zero at the
central ratio. Convergence on the boundary and the argument principle
give exactly one zero counted with multiplicity for sufficiently large
gaps, so any zero is unique and simple. Zero is an exact eigenvalue throughout;
on the simple set its projector is
\[
\Pi_p=\frac{\operatorname{adj}(-A)}{d_A}.
\]
The entire projector therefore has at most one simple pole in this
disk.

At either quadratic root, the pair formula for $L_i$ has denominator
$1+a$ and nonzero numerator. Its polarized component with all
spectators empty converges uniformly on the disk boundary to $l_+$,
with nonzero residue, because the limiting projection is simple there.
Triangularity ensures this entry uses only the diagonal grade block,
and $F=1$ there. Its Cauchy integral is consequently
nonzero for sufficiently large gaps, so $L_i$ has a pole inside the
disk. Every vector pole is a projector pole because the vectors are its
vacuum row and column. Since $\Pi_i=R_iL_i$ has at most one simple pole,
$R_i$ is analytic there: poles in both factors would give a nonzero
double-pole outer product. For $p=-i$ it is $R_{-i}$ whose pair denominator is
$1+a$; the same argument proves analyticity of $L_{-i}$.

Near $-1$, $N_2^*$ is nonzero for sufficiently separated spectators.
Its normalized determinant decouples into the pair and singleton
diagonal entries, and its pair value at $(u,-u)$ is
$-2\sqrt2\,u^2$. Away from exact shifted-site coincidences, the
determinant criterion and $\mathcal J=NN_2^*$ put every projector
pole at a zero of $N$. At the sole shift in this disk, $v=-u$, retain
the first row type of $\mathcal L$ and the sum of its two row types,
and multiply the first type for each of the two sites by $u+v$.
The resulting matrix is analytic. The first type had at most a simple
kernel pole; the sum is the rank-one inversion row, using only probes
$x_j,q^3x_j$, which meet no local pole here, and $T_p(x_j)$ is regular.
The factor $(u+v)^2$ in its scalar determinant exactly removes the
Pfaffian denominator poles. The new scalar determinant is nonzero
at the shift whenever $\mathcal J\ne0$. Thus a pole here also requires
a zero of $N$. As the projector has at most a simple pole, one factor
$N$ clears every component of each selected vector.

It remains to control zero and infinity in the varied coordinate.
Keep the other sites generic on one scale and put $\zeta$ on a separate
singleton scale. The only block carrying the character has the
singleton vacant and the other list in its generically simple fixed
character. Choose enough probes to isolate this simple zero. In the
coordinates $t=\sqrt\zeta$ near zero and $t^{-1}$ near infinity,
the original probe matrices are holomorphic by the local limits
recorded above, including the moving singleton probe $w=\zeta$.
Their simple spectral projection, and hence its vacuum row and column,
are holomorphic at the endpoint.

These limiting probes preserve the singleton's vacancy parity. If
the singleton is inactive its only possible change is by 2; if it is
active, all inactive changes are even and total charge conservation
makes its change even as well. The sign covariance then gives the
precise endpoint orders: a coordinate with $\sigma_j=\pm1$ is
$O(t)$ at zero and $O(t^{-1})$ at infinity, whereas one with
$\sigma_j=0$ is bounded. In the singleton-vacant diagonal block the
similarity $F$ is 1 and the projection is exactly the reduced
projection. Every coordinate with $\sigma_j=0$ therefore tends to
the corresponding reduced coordinate at both endpoints.

For sufficiently separated generic spectators, these arguments clear
all nonzero finite poles of $N(R_i)_\sigma$ and $N(L_{-i})_\sigma$.
Otherwise a sequence of counterexamples with growing scale gaps
would contradict the simple-probe and disk analysis. After dividing
by the displayed half powers, the endpoint orders give no pole at
zero. At infinity, $\deg_{x_j}N\le M-1$ gives degree at most
$M-1-\sigma_j^2$. The set of separated generic spectator values is
Zariski dense. A relatively prime numerator and denominator in
$x_j$ remain relatively prime for generic specializations of the
spectators, so the one-variable conclusion holds over their
coefficient field. Applying this to every $j$ removes every
irreducible denominator and proves joint polynomiality. This last
step includes simultaneous divisor intersections; the simple-pole
analysis was only needed at generic divisor points. Finally the
coordinate limits and \eqref{j:normalizer-deletion} give the stated
deletion values and top coefficients.
\end{proof}

At any nonzero site list with $N(X)\ne0$, this lemma gives
componentwise continuation of $R_i$ and $L_{-i}$ through the
polynomial formulas, including at simultaneous shifted-site
coincidences. This is the continuation used later for tilted physical
lists; the homogeneous determinant confluence above is a separate
construction.

\begin{corollary}[The scalar vacuum polynomial]\label{j:scalar-vacuum-polynomial}
For two consecutive lists $A,B$, allowing either to be empty, and
$h\ne0$, put
$\mathscr G(X;A,B)=L_-D_A R_+$. Then
\[
\begin{split}
N(X)^2\mathscr G(X;A,B)
 =\sum_{\substack{\sigma\in\{-1,0,1\}^M\\ \sum_j\sigma_j=0}}
 (ih)^{\sum_{j\in A}\sigma_j}
 \left(\prod_jx_j^{\sigma_j^2}\right)
 V_\sigma^L(X)V_\sigma^R(X).
\end{split}
\]
The empty-list value is 1. For $M\ge1$ it is a polynomial in $X$ with
Laurent polynomial coefficients in $h$, of individual degree at most
$2(M-1)$. For a deleted site $j$,
its value at $x_j=0$ is the reduced polynomial times
$\prod_{k\ne j}x_k^2$, and its coefficient of $x_j^{2(M-1)}$ is
the reduced polynomial. Moreover $\mathscr G$ itself tends to the
reduced $\mathscr G$ as $x_j\to0,\infty$ with generic spectators.
\end{corollary}
\begin{proof}
Expand the diagonal insertion in the physical spin basis and apply
Lemma~\ref{j:selected-vacuum-vectors}. A summand has degree at most
$2(M-1)-\sigma_j^2$ in $x_j$. Only $\sigma_j=0$ can contribute
to the value at zero or to the stated top coefficient, and its two
vector deletion rules give the claims. The normalized coordinate
limits in the proof of the lemma give the last assertion, also when
a group becomes empty.
\end{proof}

\subsection{Fusion of the scalar observable}

The function $\mathscr G$ is symmetric within each group. At generic
sites the Yang--Baxter similarity interchanges adjacent site
parameters, preserves charge on that pair, and works for both twists.
It therefore commutes with $D_A$ for a same-group interchange.
The fixed lines transform by this similarity; their normalization
scalars cancel in $L_-D_A R_+$ because $L_-R_+=1$ before and after
the interchange.

For a same-group pair $u,q^\delta u$, $\delta=2,3$, the function
$\mathscr G$ reduces to $\mathscr G(X_{\rm red})$. By symmetry place
$q^\delta u$ immediately before $u$. Equation~\eqref{j:site-fusion}
then intertwines the full row with a row whose site is $qu$ for
$\delta=2$, or with both sites removed for $\delta=3$. Full and
reduced fixed lines are simple at generic points of this
specialization: separate the other scales and use the active pair
test, where $a(q^\delta)=0$. Thus the specialization is taken by
regular eigenprojection. The full right vector is $J_\delta$ applied
to the reduced one, since the vacuum row of $J_\delta$ extracts the
vacuum. The full left vector restricts along $J_\delta$ to the
reduced left vector; its scalar is 1 by the overlap $L_-R_+=1$.
The insertion commutes through $J_\delta$ because the pair lies in
one group. These generic evaluations, together with
\eqref{j:normalizer-fusion}, give exact fusion evaluations of the
cleared scalar polynomial.

\subsection{Normalization on paired scales}

When the groups have equal size, the fusion and deletion evaluations
leave one scalar coefficient undetermined. We determine it by putting
one site of each group on each separated scale, with generic fixed
ratios. A charged scale block cannot carry the character. For an odd
pair total $S_j=\pm1$, the product of the two regular site probes is
$p_j^{S_j}\ne1$ by the parity test. Indeed the two transfers on active
spaces 1,2 are $R_{12}p_j^{s_1}$ and $p_j^{s_2}R_{21}$ at opposite
differences, and unitarity gives the product. For $S_j=\pm2$, the
first probe gives $p_j^{\pm1}a\ne1$ at a generic ratio.

There is consequently an isolated simple character in the product of
neutral pair blocks, with the original twists. In the vacuum matrix
element only its vacuum grade block contributes. The insertion $D_A$
is diagonal and commutes with the common diagonal similarity $F$.
That similarity therefore cancels between the two projections, and
its entry on the empty vector is 1. The result is a product of pair
projections. For a pair
whose first coordinate $u$ is in $A$ and second coordinate $v$ is
in $B$, the formulas for $l,r$ and the sine weights give
\[
\frac{d_{\rm tile}^2}{(1-a)^2}
 =(2-\sqrt2)\frac{uv}{(u+v)^2}.
\]
Thus the paired-scale limit is
\begin{equation}\label{j:paired-vacuum-limit}
\mathscr G\ \longrightarrow\ \prod_{\rm pairs}\left(1+
 (2-\sqrt2)(h+h^{-1})\frac{uv}{(u+v)^2}\right)
\end{equation}
(read at the limiting ratios).

\subsection{The scalar characterization}

Put $\mathscr F_{A,B}=N(X)^2\mathscr G(X;A,B)$, and write
$e_k(u,v)=v-q^ku$. In a deletion or fusion, the remaining coordinates
keep their groups and cyclic order; a fused site belongs to the group
of its two predecessors.

\begin{lemma}[Scalar characterization]\label{j:scalar-characterization}
The family $\mathscr F_{A,B}$ is the unique family in
$\mathbb C(h)[X]$ with the following properties. The empty-list value
is 1. For $M=|A|+|B|\ge1$, each polynomial is symmetric within the
two groups and has individual degree at most $2(M-1)$. Deleting a
site $j$ gives
\[
\left.\mathscr F_{A,B}\right|_{x_j=0}
 =\left(\prod_{k\ne j}x_k^2\right)\mathscr F_{\widehat j},
\qquad
[x_j^{2(M-1)}]\mathscr F_{A,B}=\mathscr F_{\widehat j}.
\]
For two sites $u,q^\delta u$ in one group, $\delta=2,3$, it satisfies
\[
\mathscr F_{A,B}(u,q^\delta u,Z)
 =(u+q^\delta u)^2
 \left(\prod_{z\in Z}\ell_\delta(u,z)^2\right)
 \mathscr F_{\rm red}(X_{\rm red}),
\]
with $X_{\rm red}$ as in \eqref{j:normalizer-fusion}. When
$|A|=|B|$, its quotient by $N^2$ has the paired-scale limit
\eqref{j:paired-vacuum-limit} at generic fixed pair ratios.
\end{lemma}
\begin{proof}
The preceding corollary, fusion argument and paired normalization
show that $\mathscr F_{A,B}$ has these properties. For uniqueness,
induct on $M$ and take the difference of two such families. If the
larger group has size $n>M/2$, vary one of its sites. Induction gives
zeros at the $4(n-1)$ fusion points $q^{\pm2},q^{\pm3}$ times the
other sites of that same group and at zero. Its top coefficient also vanishes,
so its degree is at most $2M-3$. Since $4(n-1)+1>2M-3$, the
difference is zero.

If the two groups have size $n$, then $M=2n$ and the vanished top
coefficient reduces the degree in each variable to at most $2M-3$.
The zero at the origin and the four zeros for each same-group partner
give degree
$1+4(n-1)=4n-3=2M-3$. Divisibility by these distinct linear
factors therefore shows that the difference is a constant in the site
variables times
\[
\prod_i x_i\prod_{\text{same-group pairs }i,j}
 e_2(x_i,x_j)e_3(x_i,x_j)e_5(x_i,x_j)e_6(x_i,x_j).
\]
To determine the constant, put matched coordinates $u_j,v_j$ on
successive increasing scales $\rho_j$. At generic ratios,
\[
N(X)\sim
\prod_j(u_j+v_j)\prod_{j<k}(u_kv_k)^2.
\]
Indeed group the pairs in the Pfaffian. Their two cross rows toward
each external pair have identical limits, so the Pfaffian factors
into its internal pair terms; $\mathcal P$ supplies the displayed
cross-scale powers. If the coefficient of the possible difference is
nonzero, its quotient by $N^2$ has nonzero limit, that coefficient times
$\prod_j u_jv_j/(u_j+v_j)^2$. The common paired-scale normalization
forces that constant to vanish.
\end{proof}

\subsection{Finite spin sum}
We now verify an explicit family with the properties of
Lemma~\ref{j:scalar-characterization}. Use indices $s,t=-1,0,1$
(symbols $-,0,+$ in tables), and let
$\kappa_{st}\in\{-1,0,1\}$ be congruent to $s-t$ modulo 3. Define
\[
p_{st}(u,v)=q^{-\kappa_{st}}P_{st},\qquad p'_{st}(u,v)=q^{\kappa_{st}}P'_{st}
\]
by the following products, listing ascending order; for descending use the swapped-index products with $k\mapsto-k$ throughout.
\[
\begin{array}{c|cc}
(st)&P&P'\\ \hline
(--),(++)&e_2 e_3 e_5 e_6&1\\
(00)&e_2 e_3 e_5 e_6/(e_1e_7)&e_4^2\\
(-0),(0+)&e_3 e_5 e_6/e_0&e_5^2\\
(-+)&e_5e_6/(e_0e_1)&(e_5 e_6)^2
\end{array}
\]
Both phased products are symmetric under simultaneous swap of spins and
arguments, since $e_k(v,u)=-q^k e_{-k}(u,v)$. Set $c=2-\sqrt2$ and
$d_{\rm a}=i\sqrt c,\ d_{\rm b}=-i\sqrt c$ on the two groups.

\begin{theorem}[Finite spin sum]\label{j:finite-spin-formula}
For $h\ne0$ and the two-list vacuum observable at generic sites,
\begin{equation}\label{j:finite-spin-sum}
\begin{aligned}
N(X)^2\mathscr G
 ={}&\sum_{\substack{(s_i)\\ \sum_{i\in A}s_i=\sum_{i\in B}s_i=Q}}
 h^Q \prod_i (d_{\sigma(i)}x_i)^{s_i^2}\\
 &\qquad\cdot\prod_{\text{same-group pairs }i,j}p_{s_i s_j}(x_i,x_j)
 \prod_{\text{cross pairs }i,j}p'_{s_i s_j}(x_i,x_j),
\end{aligned}
\end{equation}
where $\sigma(i)={\rm a},{\rm b}$ indicates group and $Q$ ranges freely
subject to the displayed equality. These auxiliary spins in group $B$
use the equal-sum convention. The combined expression is evaluated
at same-group coincidences by its polynomial limit.
\end{theorem}
\begin{proof}
We verify the properties in Lemma~\ref{j:scalar-characterization}.
Pair degrees in the same-group and cross factors are $2+2st$ and
$2-2st$. The charge constraint therefore gives termwise growth degree
at most $2(M-1)-s^2$ in a variable of spin $s$.

Poles can occur only between same-group sites. At equal sites the sum
is symmetric and has at most a simple pole, so its residue is zero.
At $v=qu$ the two singular types are $(-,+)$ and $(0,0)$. Their
spectator $P$ products agree, as do their spectator $P'$ products:
shift the indices $k$ in the second site's factors up by 1. The
phased ratios for a spectator of spin $t$ are $q^{-3t}$ in the same
group and $q^{3t}$ in the other group; their total is 1 by the
equal-sum constraint. The internal ratio of residues, including
site factors, is
\[
-c\,\frac{q-q^7}{(q-1)(q-q^2)(q-q^3)}=-1.
\]
The two residues cancel, and symmetry gives the reverse shift.
These generic divisor cancellations remove all denominators of the
sum. At zero and in the top coefficient only spin 0 survives.
Its phases against the other sites cancel by the sum constraint,
giving the reduced sum times $\prod_{j\ne i}x_j^2$ at zero and
the reduced sum in the top coefficient.

At $v=q^\delta u$ within a group the surviving terms fuse bijectively
to reduced configurations. Write their spins as $a,b$ and the
reduced spin as $\gamma$, with $*$ for a removed pair. For a spectator
spin $t$ at $w$, divide its two couplings by the reduced coupling when
there is one. The result is the monic factor $\ell_\delta(u,w)^2$
times $q^{\nu t}$ for a same-group spectator or $q^{-\nu t}$ for a
cross spectator. Divide the internal site and pair factors by the
reduced site factor, when present, and by $(u+q^\delta u)^2$; call
the result $I$. The possibilities are
\[
\begin{array}{c|cc|c|rc}
\delta&a&b&\gamma&\nu&I\\ \hline
2&-&0&-&1&q^{-1}\\
2&0&+&+&1&q\\
2&-&+&0&-2&1\\
3&-&+&*&-1&1
\end{array}
\]
All other types have a zero $e_\delta$. To check the spectator factors,
add the index multisets for $(a,t)$ and the $\delta$-shifted multiset
for $(b,t)$, counting denominator multiplicities negatively. For
$\delta=2$ subtract the 1-shifted multiset for $(\gamma,t)$. The
remaining indices are $2[5]$ or $2[5]+2[6]$, respectively, and the
$\kappa$ differences give the displayed phases. Internally,
\[
p_{-0}(1,q^2)=p_{0+}(1,q^2)=(1+q^2)^2,\quad
p_{-+}(1,q^2)=-2/c,\quad p_{-+}(1,q^3)=-1,
\]
using $(q-1)^2=-cq$. The spectator phases and $I$ cancel by the
charge constraint, and $Q$ is unchanged. This is the required fusion
multiplier.

It remains to check the paired-scale normalization. For a term with
spins $s_j,t_j$ on matched coordinates $u_j,v_j$, put
$H_j=\sum_{k\le j}(s_k-t_k)$, so $H_0=H_n=0$. Its net scaling after
division by $N^2$ is
\[
\prod_j\rho_j^{H_j^2-H_{j-1}^2}
 =\prod_{j<n}\left(\frac{\rho_j}{\rho_{j+1}}\right)^{H_j^2}.
\]
It vanishes unless every $s_j=t_j$. In that case the cross-scale
phases cancel, as do the monomials against the asymptotic for $N^2$
in the characterization proof. Inside a pair the neutral numerator
is $(u_j+v_j)^2$, and the two polarized numerators are
$c u_jv_j h^{\pm1}$. The resulting product is
\eqref{j:paired-vacuum-limit}. The empty-list sum is 1, so all properties
of Lemma~\ref{j:scalar-characterization} hold and the formula follows.
At the physical homogeneous point its quotient by $N^2$ is
$G_{A,B}$ by the homogeneous projector identification.
\end{proof}
